% To be reformatted in the journal style upon acceptance.
\documentclass[11pt]{amsart}

\usepackage{amsmath,amssymb}
\usepackage[hidelinks]{hyperref}

\newtheorem{thm}{Theorem}[section]
\newtheorem{lem}[thm]{Lemma}
\newtheorem{cor}[thm]{Corollary}
\newtheorem{prop}[thm]{Proposition}
\theoremstyle{definition}
\newtheorem{defi}[thm]{Definition}
\newtheorem{rem}[thm]{Remark}
\newtheorem{ques}[thm]{Question}
\newtheorem*{thmstar}{Theorem}

\newcommand{\DM}{\mathrm{DM}}
\newcommand{\KL}{\mathrm{KL}}
\newcommand{\lc}{Q}
\newcommand{\uc}{U}
\newcommand{\bnd}{\partial}
\newcommand{\Face}{\mathrm{Face}}
\newcommand{\rmin}{\mathrm{r}}
\newcommand{\Kar}{\mathrm{Kar}}
\newcommand{\Kthree}{\mathbf{3}}

\begin{document}

\title[The Kleene interpolation]{The Kleene interpolation:\\
  boundary fibers and the strict layer\\ across the subvarieties of
  De Morgan algebras}

\author{Ryota Shibusawa}
\address{Daiichi Institute of Technology, Japan}
\email{r-shibusawa@daiichi-koudai.ac.jp}

\subjclass[2020]{06D30; 06A07, 03B38, 55U35}
\keywords{Kleene algebra, De Morgan algebra, free algebra, cubical
  sets, boundary fiber, free spectrum, Eilenberg--Zilber category,
  cubical type theory}

\begin{abstract}
  Two companion papers determined the boundary fibers of free De
  Morgan algebras (Boolean intervals $2^{D}$ over antichains of
  diagonal points, Sperner-extremal) and the resulting structure of
  the strict layer of De Morgan cubical sets (a defect formula and a
  coherence dichotomy).  This paper computes the same theory over
  the remaining subvarieties of De Morgan algebras --- by Kalman's
  classical theorem, exactly the Boolean and Kleene varieties.  The
  free Kleene algebra on $n$ generators is identified with the
  lattice of monotone Boolean functions on the \emph{unmixed} part
  of the literal cube, the points not containing both a $(0,0)$ and
  a $(1,1)$ pair; the identification recovers the free spectra
  $6, 84, 43918, \dots$ of Berman--Mukaidono.  Only two unmixed
  points are invisible to the faces --- the poles --- and they are
  comparable, so \emph{boundary fibers in free Kleene algebras have
  at most two elements, with exactly two non-singleton fibers in
  every dimension}; realizable boundaries therefore number two fewer
  than elements, matching Berman's constant-free spectrum
  (OEIS A007154) and giving that sequence a new, cubical
  interpretation.  The defect formula and the coherence dichotomy
  port verbatim, with defects confined to the poles.  The result is
  a complete classification of strict-layer behaviour across the
  subvariety lattice: Boolean and bare distributive interval
  theories have singleton fibers (Eilenberg--Zilber-compatible
  degeneracy behaviour), Kleene has exactly two polar exceptions
  independent of dimension, and De Morgan is Sperner-extremal ---
  the three grades of involution measure the distance from
  Eilenberg--Zilber structure.  A completeness theorem closes the
  classification along the signature axis as well: nontrivial
  boundary fibers require both the involution and a connection in
  the interval signature, and once both are present De Morgan's laws
  and Kalman's theorem leave no room beyond the three grades.  All
  counts and theorems are verified exhaustively by machine through
  three generators.
\end{abstract}

\maketitle

\section{Introduction}
\label{sec:intro}

The interval theories of cubical type theory are bounded
distributive lattices with extra structure, and the free algebras of
those theories govern the \emph{strict} (judgmental) layer of the
resulting cubical sets: their elements are the reparametrizations a
proof assistant applies silently.  Over the De Morgan interval
theory of Cohen--Coquand--Huber--M\"ortberg \cite{CCHM}, two
companion papers analyzed this layer.  In \cite{DMFibers} the
\emph{boundary fibers} of the free De Morgan algebra $\DM_n$ ---
the sets of elements sharing all $2n$ face restrictions --- were
shown to be Boolean intervals $2^{D}$, with $D$ an antichain of
\emph{diagonal} points of the literal cube, of size up to
$2^{\binom{n}{\lfloor n/2 \rfloor}}$ by Sperner's theorem; in
\cite{StrictDM} this non-uniqueness was organized into a structure
theory for the De Morgan cube category (a one-step \emph{defect
formula} and an exact \emph{coherence dichotomy}), against the
background fact that the category is not Eilenberg--Zilber.

By a classical theorem of Kalman \cite{Kalman}, the variety of De
Morgan algebras has exactly two proper nontrivial subvarieties:
Boolean algebras and \emph{Kleene algebras} (De Morgan algebras
satisfying $x \wedge \neg x \leq y \vee \neg y$) --- the three-element
chain $\Kthree$ being the generating subdirectly irreducible of the
latter.  Each subvariety has its cube category and its cubical sets
\cite{BuchholtzMorehouse}, and it is natural to ask how the fiber
and strict-layer theory varies across the lattice of subvarieties.
This paper answers the question completely; the summary is the
following table, in which the last row is \cite{DMFibers,StrictDM}
and the first is folklore.

\begin{center}
\small
\begin{tabular}{l|l|l|l}
  theory & invisible pts & fibers &
  non-singleton fibers\\
  \hline
  Dedekind, Boolean & none & $1$ & none\\
  \textbf{Kleene} & the two poles & $\leq 2$ &
  \textbf{exactly $2$, every $n$}\\
  De Morgan & $2^n$ diagonals & $\leq 2^{\binom{n}{\lfloor
  n/2\rfloor}}$ & $34.5\%$ of $\DM_3$
\end{tabular}
\end{center}

\noindent The Kleene row is the content of this paper, and it sits
strictly --- and, as it turns out, very asymmetrically --- between
its neighbours: the Kleene law does not restore unique fillers, but
it cuts the exceptional locus from a positive fraction of all
elements down to \emph{exactly two fibers, independent of
dimension}.

The technical basis is a representation theorem
(Theorem~\ref{thm:rep}) that seems worth recording in its own
right.  Call a point of the literal cube $\{0,1\}^{2n}$
\emph{mixed} if some coordinate pair equals $(0,0)$ and another
equals $(1,1)$, and let $\uc_n$ be the subposet of unmixed points.
Then the free Kleene algebra $\KL_n$ is the lattice of monotone
Boolean functions on $\uc_n$ --- the Kleene law deletes from the De
Morgan representation \cite{DMFibers} exactly the mixed points.
The proof is a two-sided probe of the three-valued semantics: the
lower and upper Boolean approximations of a point of $\Kthree^n$
are always unmixed, and every unmixed point arises as such an
approximation.  The cardinalities
$|\KL_n| = 6, 84, 43918, \dots$ then agree with the free spectra
computed by Berman--Mukaidono \cite{BermanMukaidono}
(OEIS~A268534), which we take as independent confirmation.

The fiber theory follows in one stroke
(Theorem~\ref{thm:fibers}): the points of $\uc_n$ invisible to all
faces are the two poles $\bot, \top$ --- every other diagonal point
of the De Morgan literal cube is mixed --- and the poles are
comparable, so the defect antichains of \cite{DMFibers,StrictDM}
can only be $\emptyset$, $\{\bot\}$ or $\{\top\}$.  Fibers have at
most two elements; exactly two fibers are non-singleton (the
\emph{polar} fibers $\{0, \delta_{\top}\}$ and
$\{1 \setminus \bot, 1\}$); and the number of realizable boundaries
is exactly $|\KL_n| - 2$.  Numerically this is the sequence
$4, 82, 43916, \dots$ --- which turns out to be in the OEIS as
A007154, \emph{Berman's constant-free free spectrum of Kleene
algebras} \cite{Berman}; the coincidence is explained by a
two-line bijection (boundaries $\leftrightarrow$ fiber-minimal
elements, of which all but two elements are), giving the classical
sequence a new interpretation as a count of spheres in the Kleene
cube category.

Finally, the strict-layer theory of \cite{StrictDM} ports verbatim
(\S\ref{sec:strict}): the Kleene cube category is not idempotent
complete (the proof for De Morgan applies word for word, as
$|\KL_1| = |\DM_1| = 6$), its Karoubi envelope is the opposite of
finite projective Kleene algebras via Bova--Cabrer
\cite[Thm.~12]{BovaCabrer}, the defect formula holds with defects
confined to the poles, and the coherence dichotomy is exact --- with
the exceptional locus of any substitution now consisting of at most
the two polar elements.  All statements are verified exhaustively
by machine for $n \leq 3$ (\S\ref{sec:machine}).

\subparagraph*{Reading guide.}  This paper is the third of a
series but is self-contained in its statements; proofs that are
verbatim ports from \cite{DMFibers,StrictDM} are indicated as such,
with the one substitution ``diagonal points $\mapsto$ poles''
spelled out.  The conceptual through-line is stated once, here: the
lattice of subvarieties of De Morgan algebras (trivial $\subset$
Boolean $\subset$ Kleene $\subset$ De Morgan, by Kalman) is
simultaneously a scale of \emph{how badly the associated cube
category fails to be Eilenberg--Zilber}, with the failure measured
by boundary-fiber structure --- none, two polar bits, a full
Sperner antichain of bits.  Section~\ref{sec:complete} shows the
scale is complete: no interval signature produces any other fiber
behaviour.

\section{The representation theorem}
\label{sec:rep}

A \emph{Kleene algebra} is a De Morgan algebra satisfying
$x \wedge \neg x \leq y \vee \neg y$; the variety is generated by
the three-element chain $\Kthree = \{0, u, 1\}$ with $\neg u = u$
\cite{Kalman, GWW}.  Write $\KL_n$ for the free Kleene algebra on
$x_1, \dots, x_n$.  We use the literal-cube representation of
$\DM_n$ from \cite{DMFibers}: monotone functions (up-sets) on
$\lc_n = \{0,1\}^{2n}$, coordinates in pairs $(x_i, \neg x_i)$.

\begin{defi}
  A point $\alpha \in \lc_n$ is \emph{mixed} if some pair
  $\alpha_i = (0,0)$ and some pair $\alpha_j = (1,1)$; otherwise
  \emph{unmixed}.  Let $\uc_n \subseteq \lc_n$ be the subposet of
  unmixed points.  ($|\uc_1| = 4$, $|\uc_2| = 14$,
  $|\uc_3| = 46$.)
\end{defi}

Note that $\uc_n$ is closed under the De Morgan symmetry
(complement all coordinates, swap each pair), and that the
\emph{diagonal} points of $\lc_n$ (all pairs in
$\{(0,0), (1,1)\}$) that are unmixed are exactly the two poles
$\bot = (0,\dots,0)$ and $\top = (1,\dots,1)$.

\begin{thm}
  \label{thm:rep}
  Restriction of monotone functions along
  $\uc_n \hookrightarrow \lc_n$ induces an isomorphism from the
  free Kleene algebra onto the lattice of monotone functions on
  $\uc_n$, with involution
  $(\neg\varphi)(\alpha) = 1 - \varphi(\sigma\bar\alpha)$ as
  before:
  \[
    \KL_n \;\cong\; \{\text{monotone } \uc_n \to \{0,1\}\}.
  \]
\end{thm}

\begin{proof}
  The restriction map
  $\rho : \DM_n \to \mathrm{Mon}(\uc_n)$ is a homomorphism of De
  Morgan algebras: meets, joins and bounds are pointwise, and
  $\uc_n$ is closed under the symmetry defining $\neg$.  It is
  surjective, since a monotone function on $\uc_n$ extends to
  $\lc_n$ (e.g.\ by its least monotone extension).

  $\mathrm{Mon}(\uc_n)$ is a Kleene algebra: if
  $\alpha \in x_i \wedge \neg x_i$ then $\alpha_i = (1,1)$, so by
  unmixedness no pair of $\alpha$ is $(0,0)$; hence
  $\alpha_j \in \{(0,1), (1,0), (1,1)\}$ and
  $\alpha \in x_j \vee \neg x_j$, for every $j$.  Thus $\rho$
  factors through the Kleene quotient, giving a surjection
  $\KL_n \twoheadrightarrow \mathrm{Mon}(\uc_n)$.

  For injectivity we use the three-valued semantics.  By
  \cite{GWW}, $\KL_n$ embeds into $\Kthree^{(\Kthree^n)}$ as the
  subalgebra generated by the projections; write
  $\hat\varphi : \Kthree^n \to \Kthree$ for the function induced by
  $\varphi$.  For $v \in \Kthree^n$ define
  $v_{\mathrm{low}}, v_{\mathrm{high}} \in \lc_n$ by encoding
  $v_i = 0 \mapsto (0,1)$, $v_i = 1 \mapsto (1,0)$, and
  $v_i = u \mapsto (0,0)$ (low), resp.\ $(1,1)$ (high).  Both are
  \emph{unmixed}, because the choice for $u$ is made uniformly; and
  a direct induction on terms shows
  $\hat\varphi(v) = 0, u, 1$ according as the pair
  $\bigl(\varphi(v_{\mathrm{low}}), \varphi(v_{\mathrm{high}})\bigr)$
  is $(0,0)$, $(0,1)$, $(1,1)$ (monotonicity excludes $(1,0)$).
  Hence $\hat\varphi$ is determined by $\rho(\varphi)$.
  Conversely every unmixed point is a $v_{\mathrm{low}}$ or a
  $v_{\mathrm{high}}$ (read $u$ at its diagonal pairs), so
  $\rho(\varphi)$ is determined by $\hat\varphi$; as
  $\varphi \mapsto \hat\varphi$ is injective on $\KL_n$, so is the
  induced map $\KL_n \to \mathrm{Mon}(\uc_n)$.
\end{proof}

\begin{rem}
  The counts $|\mathrm{Mon}(\uc_n)| = 6, 84, 43918$ for
  $n = 1, 2, 3$ (machine enumeration, \S\ref{sec:machine}) agree
  with the free spectra of Kleene algebras computed by
  Berman--Mukaidono \cite{BermanMukaidono} (OEIS~A268534), an
  independent confirmation of Theorem~\ref{thm:rep}.  The theorem
  refines the classical normal-form descriptions \cite{GWW,
  BermanMukaidono} in the same way the literal-cube picture refines
  them for De Morgan algebras: it makes the \emph{boundary}
  combinatorics visible.  In that picture the Kleene law acts by
  \emph{deleting the mixed points} --- and, as the next section
  shows, mixed points are almost all of the diagonal.
\end{rem}

\section{Boundary fibers: at most two elements}
\label{sec:fibers}

Faces $\bnd_i^e : \KL_n \to \KL_{n-1}$ restrict to the face
subcubes of $\uc_n$ exactly as in the De Morgan case (the face
subcubes of $\uc_n$ are copies of $\uc_{n-1}$: fixing a pair to
$(0,1)$ or $(1,0)$ cannot create or destroy mixedness).  A point of
$\uc_n$ is a \emph{face point} if some pair is $(0,1)$ or $(1,0)$;
the points invisible to every face are the unmixed diagonal
points:

\begin{lem}
  \label{lem:poles}
  The unmixed diagonal points of $\lc_n$ are exactly the poles
  $\bot$ and $\top$; they are comparable, and there is a face point
  strictly between them (any point with one pair $(0,1)$ and the
  rest $(0,0)$).
\end{lem}

\begin{thm}
  \label{thm:fibers}
  Every boundary fiber of $\KL_n$ is an interval $[m, M]$ with
  $D = M \setminus m$ an antichain contained in
  $\{\bot, \top\}$; hence $|D| \leq 1$ and fibers have at most two
  elements.  Exactly two fibers are non-singleton in every
  dimension $n \geq 1$: the \emph{zero-boundary} fiber
  $\{0,\ \delta_{\top}\}$ ($\delta_{\top}$ the indicator of
  $\top$, i.e.\ $x_1 \wedge \neg x_1 \wedge \cdots \wedge x_n
  \wedge \neg x_n$) and the \emph{top-boundary} fiber
  $\{1 \setminus \bot,\ 1\}$.  Consequently
  \[
    B_{\KL}(n) \;=\; |\KL_n| - 2 ,
  \]
  where $B_{\KL}(n)$ is the number of realizable boundaries.
\end{thm}

\begin{proof}
  The interval property, the confinement of $D$ to invisible
  points, and the antichain property are the proofs of
  \cite[\S3]{DMFibers} verbatim, applied to the poset $\uc_n$: the
  only inputs are that faces see exactly the face points and that
  between comparable invisible points there is a face point
  (Lemma~\ref{lem:poles}).  Since $\bot < \top$, an antichain
  contains at most one pole.  Realizability of $D = \{\top\}$ and
  $D = \{\bot\}$ is witnessed by the displayed fibers (that
  $\delta_{\top}$ is monotone on $\uc_n$ and has zero boundary, and
  dually, is immediate); uniqueness of each --- i.e.\ that no other
  fiber has $D \neq \emptyset$ --- follows from the fiber
  description of \cite[Thm.~3.4]{DMFibers}: $D = \{\top\}$ forces
  every face point above $\top$ to lie in $m$, and there are none,
  so $m$ is the least up-set with empty intersection with
  $\{\top\}$'s up-set, namely $m = 0$; dually for $\bot$, $m$ must
  contain every face point, forcing $m = 1 \setminus \bot$.
  Finally, boundaries biject with fiber-minimal elements
  \cite[Thm.~4.2]{DMFibers}, and all elements are fiber-minimal
  except the top members $\delta_{\top}$ and $1$ of the two polar
  fibers.
\end{proof}

\begin{rem}[the free spectrum, interpreted]
  \label{rem:spectrum}
  The sequence $B_{\KL}(n) = 4, 82, 43916, \dots$ appears in the
  OEIS as A007154, \emph{the free spectrum of Kleene algebras in
  the constant-free signature} (Berman \cite{Berman}): dropping the
  constants $0, 1$ from the signature removes exactly two elements
  from each free algebra.  Theorem~\ref{thm:fibers} gives the same
  numbers a different meaning --- the number of compatible spheres
  in the Kleene cube category --- via the bijection
  boundaries $\leftrightarrow$ fiber-minimal elements.  We do not
  know a direct structural link between the two interpretations
  beyond the count; both ``remove two elements'', but different
  ones ($\{0,1\}$ versus $\{\delta_\top, 1\}$).
\end{rem}

\begin{rem}[Boolean and Dedekind]
  \label{rem:boolean}
  For the Boolean cube category ($x \vee \neg x = 1$), the
  representation collapses to functions on the \emph{consistent}
  points (all pairs $(0,1)$ or $(1,0)$): every point is a face
  point, no invisible points remain, and all fibers are singletons
  --- boundaries determine fillers, as for the bare distributive
  (Dedekind) theory \cite{DMFibers}.  With Kalman's theorem
  \cite{Kalman} --- Boolean and Kleene are the only proper
  nontrivial subvarieties of De Morgan algebras --- the table of
  \S\ref{sec:intro} is therefore a \emph{complete} classification
  of boundary-fiber behaviour over the subvariety lattice.
\end{rem}

\section{The strict layer}
\label{sec:strict}

We record how the structure theory of \cite{StrictDM} specializes.
Throughout, $\rmin(g) = g \setminus A_g$ where $A_g$ is the set of
minimal invisible points of $g$; by Theorem~\ref{thm:fibers},
$A_g \subseteq \{\bot, \top\}$, and indeed
$A_g = \{\top\}$ iff $g = \delta_\top$, $A_g = \{\bot\}$ iff
$g = 1$.

\begin{prop}
  \label{prop:kar}
  The Kleene cube category is not idempotent complete
  ($x \mapsto x \vee \neg x$ does not split; the proof of
  \cite[Prop.~3.1]{StrictDM} applies verbatim since
  $|\KL_1| = 6$), hence not Eilenberg--Zilber; and its Karoubi
  envelope is the opposite of the category of finite projective
  Kleene algebras, characterized combinatorially by Bova--Cabrer
  \cite[Thm.~12]{BovaCabrer}.
\end{prop}

\begin{thm}
  \label{thm:strict}
  Over the Kleene cube category:
  \begin{enumerate}
  \item (\emph{defect formula}) for every $g$, $\rmin(g)$ is the
    least element of the boundary fiber of $g$ in one deletion
    step;
  \item (\emph{coherence dichotomy}) for every substitution $f$
    (inducing $F$ on unmixed points --- substitutions preserve
    unmixedness) and every $g$:
    $\rmin(f^{*}g) = \rmin(f^{*}\rmin(g))$ iff
    $F(\Face(\uc_m)) \cap A_g = \emptyset$;
  \item (\emph{polar exceptional locus}) consequently the
    incoherent pairs $(f, g)$ have
    $g \in \{\delta_{\top}, 1\}$: for each substitution, at most
    two elements --- independent of $n$ --- fail to transport
    coherently, in contrast with the positive-fraction exceptional
    loci of the De Morgan case \cite{StrictDM}.
  \end{enumerate}
\end{thm}

\begin{proof}
  (1) and (2) are the proofs of \cite[Thm.~4.1, Thm.~5.2]{StrictDM}
  verbatim over the poset $\uc_n$, with
  Lemma~\ref{lem:poles} supplying the two ingredients used there
  (faces see exactly face points; a face point lies between
  comparable invisible points).  Preservation of unmixedness by
  substitutions is the semantic fact that lower/upper
  approximations are computed uniformly:
  $F(v_{\mathrm{low}}) = (t(v))_{\mathrm{low}}$ for the
  three-valued evaluation $t$ of the substitution.  (3) is
  immediate from $A_g \subseteq \{\bot,\top\}$ and the
  identification of the two elements with nonempty defect.
\end{proof}

\begin{rem}
  Thus the Kleene strict layer is ``Eilenberg--Zilber up to two
  poles'': the minimal-filler retraction is a strictly natural
  normal form away from $\{\delta_\top, 1\}$, and even there the
  deviation is a single toggled pole, corrected inside a two-element
  fiber by the median cylinder (whose Kleene instance is the same
  term as in \cite{DMFibers}).  It is instructive that the two
  exceptional elements are precisely the witnesses of the
  \emph{failure of excluded middle} surviving the Kleene law:
  $\delta_\top = \bigwedge_i (x_i \wedge \neg x_i) \neq 0$ and its
  dual.  The Kleene law tames, but cannot eliminate, the no-LEM
  cells --- and they are exactly what separates the Kleene cube
  category from Eilenberg--Zilber behaviour.
\end{rem}

\section{Completeness: fiber behaviours across signatures}
\label{sec:complete}

The trichotomy of the introduction classifies fiber behaviour over
the subvarieties of De Morgan algebras.  One may still ask whether
\emph{weaker signatures} --- interval theories that drop the
involution or the connections --- realize intermediate behaviours.
They do not.

\begin{lem}
  \label{lem:signature}
  An interval signature (a subset of $\{\wedge, \vee, \neg\}$
  over the bounds $0, 1$) containing $\neg$ and at least one
  connection defines the other connection by De Morgan's laws.
  Hence the signature classes are: $\emptyset$, $\{\wedge\}$,
  $\{\vee\}$, $\{\wedge,\vee\}$ (the Dedekind theory),
  $\{\neg\}$, and the full De Morgan signature.
\end{lem}

\begin{thm}[completeness]
  \label{thm:complete}
  Over every interval signature not containing both the involution
  and a connection, all boundary fibers of the free algebras are
  singletons.  Consequently the realizable fiber behaviours of
  interval theories are exactly three: trivial (Dedekind, Boolean,
  and all weaker signatures), two polar fibers (Kleene), and
  Sperner-extremal (De Morgan).
\end{thm}

\begin{proof}
  By Lemma~\ref{lem:signature} four signature classes remain below
  the full one.  \emph{Dedekind}: the free bounded distributive
  lattice on $n$ generators is the lattice of monotone functions on
  $\{0,1\}^n$; every evaluation point has a first coordinate in
  $\{0,1\}$, so an element is determined by its two faces
  $x_1 := 0, 1$ --- fibers are singletons.  (The same two-face
  decomposition handles the \emph{Boolean} algebra of all
  functions, re-proving the first line of the table.)
  \emph{Meet-only} (dually join-only): the free bounded
  $\wedge$-semilattice consists of $0$ and the conjunctions
  $\bigwedge_{i \in S} x_i$; the faces of a conjunction determine
  $S$ (whether $i \in S$ is read off the $x_i := 0$ face), and
  distinguish $0$ from $1$.  \emph{Involution-only}: the free
  algebra is $\{0, 1, x_1, \neg x_1, \dots\}$, and the faces of
  $x_i$, $\neg x_i$ (constants $0, 1$ in opposite positions)
  separate all elements.  \emph{Bare}: immediate.  In each case no
  evaluation point is invisible to the faces; the nontrivial
  behaviours require literals to decouple ($\neg$, giving
  independent pairs) \emph{and} to combine ($\wedge$ or $\vee$),
  and once both are available the algebra is a De Morgan algebra,
  where Kalman's theorem \cite{Kalman} and the fiber theorems
  (\cite{DMFibers}, Theorem~\ref{thm:fibers}) leave exactly the
  three grades.
\end{proof}

\begin{rem}
  The moral sharpens the through-line of the introduction:
  \emph{fiber freedom is born only from the interaction of the
  involution with a connection, and once born it is exhausted by
  Kalman's three grades}.  In particular no interval \emph{theory}
  interpolates strictly between the Kleene and De Morgan rows ---
  any graded interpolation must come from non-equational structure
  (restricted substitution classes, cf.\ the coherence analysis of
  \cite{StrictDM}), not from the algebra.
\end{rem}

\section{Machine verification}
\label{sec:machine}

All counts and theorems were verified exhaustively for
$n \leq 3$ (scripts accompany the artifact):
$|\mathrm{Mon}(\uc_n)| = 6, 84, 43918$ against the
Berman--Mukaidono spectra; fiber censuses
$\{1^2, 2^2\}$, $\{1^{80}, 2^{2}\}$, $\{1^{43914}, 2^{2}\}$
confirming Theorem~\ref{thm:fibers} and the identity
$B_{\KL}(n) = |\KL_n| - 2$; the defect formula on all $84$
elements of $\KL_2$ against brute-force fiber minima; and the
coherence dichotomy on all of $\KL_2$ for five generating
substitutions --- conditions of Theorem~\ref{thm:strict}(2)
coincide exactly, with incoherent loci of sizes $2, 2, 0, 2, 0$
(both connections, reversal, diagonal,
$x \mapsto x \wedge \neg x$), each consisting of the two polar
elements.  Assertions that substitutions preserve unmixedness and
that defects stay within $\{\bot, \top\}$ were checked at every
step.  For Theorem~\ref{thm:complete}, the
singleton-fiber claims were enumerated as well: Dedekind
$n \leq 3$ ($6$ and $20$ elements), meet-only $n \leq 3$ ($5$ and
$9$ elements), involution-only $n = 2$ ($6$ elements) --- all
fibers singleton.  At $n = 2$ the boundary counts of the Kleene and De Morgan
theories coincide ($82$): the mixed points of $\lc_2$ are the two
non-polar diagonals, so the two face-point posets agree ---
a low-dimensional accident that disappears at $n = 3$
($43{,}916$ versus $4{,}829{,}136$ \cite{DMFibers}).

\section{Outlook}
\label{sec:outlook}

Two questions remain.  First, with the equational axis closed by
Theorem~\ref{thm:complete}, graded interpolations between the
Kleene and De Morgan rows can only arise from \emph{non-equational}
data --- wide subcategories of the cube category (restricting which
substitutions act) or filtered mixing constraints; there the
varying quantity is not the fiber but the \emph{coherence}
behaviour of \cite{StrictDM}, and a classification over the
generator classes of \cite{BuchholtzMorehouse} seems within
reach.  Second, \emph{the strict--weak interface}: in the De Morgan case
the diagonal points simultaneously index fiber freedom, coherence
defects, and the failure of idempotent splitting; the Kleene case
shows the first two can be cut to a constant while the third
persists ($x \vee \neg x$ still does not split).  Whether the
Karoubi envelope of the Kleene cube category — the finite projective
Kleene algebras of \cite{BovaCabrer} — admits a boundary-fiber
theory of its own, and how the polar cells behave there, is open.

\subparagraph*{Artifact.}  Scripts and machine-checked companions
are available in the repository at
\url{https://github.com/r-shibusawa/cubical-strict-j}
(archived at DOI \href{https://doi.org/10.5281/zenodo.21405961}%
{10.5281/zenodo.21405961}).

\begin{thebibliography}{10}

\bibitem{Berman}
J.~Berman.
\newblock Free spectra of $3$-element algebras.
\newblock In \emph{Universal Algebra and Lattice Theory}
  (R.~S.~Freese, O.~C.~Garcia, eds.), Lecture Notes in
  Math.\ 1004, Springer, 1983.

\bibitem{BermanMukaidono}
J.~Berman, M.~Mukaidono.
\newblock Enumerating fuzzy switching functions and free Kleene
  algebras.
\newblock \emph{Comput.\ Math.\ Appl.}\ 10 (1984), 25--35.

\bibitem{BovaCabrer}
S.~Bova, L.~Cabrer.
\newblock Unification and projectivity in De Morgan and Kleene
  algebras.
\newblock \emph{Order}\ 31 (2014), 159--187.

\bibitem{BuchholtzMorehouse}
U.~Buchholtz, E.~Morehouse.
\newblock Varieties of cubical sets.
\newblock \emph{RAMiCS 2017}, LNCS 10226, 2017.

\bibitem{CCHM}
C.~Cohen, T.~Coquand, S.~Huber, A.~M\"ortberg.
\newblock Cubical type theory: a constructive interpretation of the
  univalence axiom.
\newblock \emph{TYPES 2015}, LIPIcs 69, 2018.

\bibitem{GWW}
M.~Gehrke, C.~Walker, E.~Walker.
\newblock Normal forms and truth tables for fuzzy logics.
\newblock \emph{Fuzzy Sets and Systems}\ 138 (2003), 25--51.

\bibitem{Kalman}
J.~A.~Kalman.
\newblock Lattices with involution.
\newblock \emph{Trans.\ Amer.\ Math.\ Soc.}\ 87 (1958), 485--491.

\bibitem{DMFibers}
R.~Shibusawa.
\newblock Boundary fibers in free De Morgan algebras: Boolean
  intervals, Sperner bounds, and median contractions.
\newblock Preprint, 2026.  Included in the artifact repository
  \url{https://github.com/r-shibusawa/cubical-strict-j} (from
  release \texttt{v1.6.0}, DOI
  \href{https://doi.org/10.5281/zenodo.21695307}%
  {10.5281/zenodo.21695307}).

\bibitem{StrictDM}
R.~Shibusawa.
\newblock The strict layer of De Morgan cubical sets: minimal
  fillers, a coherence dichotomy, and the Karoubi envelope.
\newblock Preprint, 2026.  Included in the artifact repository
  above (from release \texttt{v1.7.0}, DOI
  \href{https://doi.org/10.5281/zenodo.21695641}%
  {10.5281/zenodo.21695641}).

\end{thebibliography}

\end{document}
